% !TeX root = ../../Book2.tex
\section{光滑依赖与流}
\label{b2:sec:ode-smooth-flow}

唯一性使从不同时间出发的轨道能够相互接合，光滑依赖则允许把轨道用于坐标变换。本节在自治 \(C^1\) 条件下建立这两项性质；高阶光滑依赖尚待撰写。

\subsection{流映射与复合关系}
\label{b2:subsec:070202}

设 \(U\subset\mathbb R^d\) 开，\(b\in C^1(U;\mathbb R^d)\)、\(a\in U\)。按\secref{b2:sec:ode-continuation}的拼接构造，初值问题 \(X'=b(X)\)、\(X(0)=a\) 有最大解区间 \(I_a\)。把从 \(a\) 出发的最大解记为 \(\Phi_t(a)\)，称 \(a\mapsto\Phi_t(a)\) 为\term[liu-yingshe]{流映射}[flow map]。记号 \(\Phi_t(a)\) 只在 \(t\in I_a\) 时使用；给映射命名并不意味着它对所有时间存在。
\symidx{\(\Phi_t\)：自治速度场的流映射}

\begin{proposition}\label{b2:prop:07-flow-composition}
设 \(d\) 为正整数，\(U\subset\mathbb R^d\) 为开集，\(b\in C^1(U;\mathbb R^d)\)，\(\Phi_t(a)\) 是 \(X'=b(X)\)、\(X(0)=a\) 的最大局部流。凡各项有定义，流满足
\[
 \Phi_0(a)=a,\quad
 \Phi_{t+s}(a)=\Phi_t\bigl(\Phi_s(a)\bigr),\quad
 \Phi_{-t}\bigl(\Phi_t(a)\bigr)=a.
\]
固定 \(t\) 时，\(\Phi_t\) 在其定义域上单射，其逆为相应的 \(\Phi_{-t}\)。
\end{proposition}
\begin{proof}
曲线 \(r\mapsto\Phi_{s+r}(a)\) 与 \(r\mapsto\Phi_r(\Phi_s(a))\) 在 \(r=0\) 具有相同初值，且都满足自治方程。唯一性给出复合关系。令 \(r=-s\) 得逆关系，进而得到单射性。
\end{proof}

\subsection{一阶变分方程}
\label{b2:subsec:ode-first-variation}

\begin{theorem}\label{b2:thm:07-flow-derivative}
设 \(d\) 为正整数，\(U\subset\mathbb R^d\) 为开集，\(b\in C^1(U;\mathbb R^d)\)，\(\Phi_t(a)\) 是 \(X'=b(X)\)、\(X(0)=a\) 的最大局部流。流 \((t,a)\mapsto\Phi_t(a)\) 在其开定义域上属于 \(C^1\)。对初始点的导数 \(A(t,a)=D_a\Phi_t(a)\) 满足
\begin{equation}\label{b2:eq:07-variational-ode}
 A'(t,a)=Db\bigl(\Phi_t(a)\bigr)A(t,a),\quad A(0,a)=I.
\end{equation}
因此固定时间的流是其定义域到像集的 \(C^1\) 微分同胚。
\end{theorem}
\begin{proof}
先在\cref{b2:thm:07-local-ode}构造的紧球和短时间内工作。式\eqref{b2:eq:07-flow-lipschitz}与 \(|\partial_t\Phi_t|\leq M\) 保证流联合连续。矩阵积分方程
\[
 A(t)=I+\int_0^t Db\bigl(\Phi_s(a)\bigr)A(s)\,\mathrm ds
\]
的系数连续，逐列应用\cref{b2:prop:ode-continuous-linear}即可构造解；其范数的积分不等式与\cref{b2:lem:ode-gronwall}给出 \(\|A(t)\|\leq e^{L|t|}\)。现取小向量 \(v\)，令 \(\delta_v(t)=\Phi_t(a+v)-\Phi_t(a)\)。中值积分给出
\[
 b\bigl(\Phi_t(a+v)\bigr)-b\bigl(\Phi_t(a)\bigr)
 =Db\bigl(\Phi_t(a)\bigr)\delta_v(t)+r_v(t).
\]
导数 \(Db\) 在紧球上一致连续。记其连续模为 \(\omega(\eta)\)，即距离至多 \(\eta\) 的两点上导数差的上确界；则 \(\omega(\eta)\to0\)。由 \(|\delta_v(t)|\leq e^{Lh}\cdot|v|\)，中值积分的余项满足
\[
 \sup_{|t|\leq h}|r_v(t)|
 \leq e^{Lh}\cdot|v|\cdot\omega(e^{Lh}\cdot|v|).
\]
令 \(e_v=\delta_v-Av\)，相减后的积分方程给出
\[
 |e_v(t)|\leq L\int_0^{|t|}|e_v(\pm s)|\,\mathrm ds
              +h\sup_{|s|\leq h}|r_v(s)|,
\]
其中正负号与 \(t\) 的符号一致。Gronwall 估计因此给出
\[
 \sup_{|t|\leq h}|e_v(t)|
 \leq h e^{2Lh}\cdot|v|\cdot\omega(e^{Lh}\cdot|v|).
\]
除以 \(|v|\) 并令 \(v\to0\)，得到
\[
 \sup_{|t|\leq h}
 \frac{|\Phi_t(a+v)-\Phi_t(a)-A(t,a)v|}{|v|}\longrightarrow0.
\]
这证明的是完整 Fréchet 微分，而非只沿坐标求导。若 \(a\) 改为邻近点，矩阵方程的系数一致收敛；作差和 Gronwall 估计又给出 \(A\) 对 \((t,a)\) 连续。时间导数 \(\partial_t\Phi_t(a)=b(\Phi_t(a))\) 也连续，因此流为 \(C^1\)。

对一段已存在的轨道，其图像在任意紧时间区间内为紧集。用上述局部构造覆盖它，并沿轨道取有限分段，连续依赖使邻近初始点的解也能逐段通过这些邻域。故流的定义域为开集，局部可微结论可以沿轨道继续。最后，\cref{b2:prop:07-flow-composition}给出双方均 \(C^1\) 的逆映射，所以导数矩阵可逆，像集也为开集。
\end{proof}

流的空间可微性比轨道本身可微更强。前者控制相邻粒子之间的距离、角度和体积，后者只说明单个粒子的速度存在。后面特征线法中 Jacobi 行列式的演化式依赖的是\eqref{b2:eq:07-variational-ode}，不能只从每条轨道 \(C^1\) 便直接写下。

\subsection{高阶光滑依赖}
\label{b2:subsec:ode-higher-dependence}

待撰写。本节目前只证明自治 \(C^1\) 速度场的一阶光滑依赖。
